% Back cover (4e de couverture): English and French abstracts on a single even page
\clearpage
\ifodd\value{page}\hbox{}\thispagestyle{empty}\clearpage\fi
\thispagestyle{empty}
\enlargethispage{3\baselineskip}
\vspace*{-1.2cm}
\begin{center}
    {\large \textbf{Optical control of ultracold $^{87}$Sr, from one-body spin manipulations to two-body photoassociation}}\\[0.4cm]
    {\large \textbf{Abstract}}\\[-0.3cm]
\end{center}

{\footnotesize
Strontium 87 is a fermionic alkaline-earth atom whose large nuclear spin ($I=9/2$) is
decoupled from the electronic shell, which gives a ten-dimensional spin
space. This thesis exploits this spin space in two directions: the coherent control of the
nuclear spin, and the photoassociation of $^{87}$Sr near the narrow intercombination line.

After describing the preparation of a degenerate Fermi gas of $^{87}$Sr, we demonstrate the
coherent manipulation of its nuclear spin. Using a tensor light shift to isolate chosen
subsystems of Zeeman sublevels, we effectively produce qubits and qudits. Ramsey
interferometry shows that the qubit coherence is preserved over $3$~s, without magnetic shielding, thanks to
the very small nuclear Land\'e factor. Building on this platform, we measure two physical
parameters of the atomic environment in a single experimental cycle, and two non-commuting observables of a qubit in a
single realization, by mapping it onto ancillary spin states; we describe this last scheme as a
positive-operator-valued measure (POVM).

We then study photoassociation near the $^1S_0\rightarrow{}^3P_1$ line. Lines of the bosonic
isotope $^{88}$Sr are used to validate the methods that extract the optical length from the
atom losses. For $^{87}$Sr, up to a few
thousand narrow lines are expected from theory within the explored range of 400~MHz. We observe only three lines,
broadened to several MHz, two to three orders of magnitude above their natural width, with
optical lengths one to two orders of magnitude below those of $^{88}$Sr. We finally discuss the possible 
interpretations of our results, chief among them the mechanism of predissociation, quantum statistics, and spin selection effects.

\vspace{0.15cm}
\noindent\textbf{Keywords:} strontium 87, nuclear spin, qudits, SU(N) symmetry, Ramsey interferometry, POVM, photoassociation, degenerate Fermi gas.
\par}

\vspace{0.1cm}
\noindent\rule{\textwidth}{0.4pt}
\vspace{0.3cm}

\begin{center}
    {\large \textbf{Contrôle optique du $^{87}$Sr ultrafroid: des manipulations de spin à un corps à la photoassociation de deux corps}}\\[0.4cm]
    {\large \textbf{Résumé}}\\[-0.3cm]
\end{center}

{\footnotesize
Le strontium 87 est un atome alcalino-terreux fermionique dont le grand spin nucléaire
($I=9/2$) est découplé du nuage électronique, ce qui donne un espace de Hilbert de dimension
dix. Cette thèse comporte des résultats qui s'appuient sur l'utilisation de ce large espace
selon deux axes~: le contrôle cohérent du spin nucléaire, et la photoassociation du $^{87}$Sr
sous la raie d'intercombinaison.

Après avoir détaillé la préparation d'un gaz de Fermi dégénéré de $^{87}$Sr, nous démontrons
la manipulation cohérente de son spin nucléaire. En utilisant un déplacement lumineux
tensoriel pour isoler des sous-systèmes de l'ensemble des sous-niveaux Zeeman, nous créons
effectivement des qubits et des qudits. L'interféromètre de Ramsey ainsi montre que
la cohérence des qubits est préservée sur $3$~s, sans isolation magnétique, grâce au très
faible facteur de Landé nucléaire. À partir de cette plateforme, nous mesurons dans deux
expériences différentes deux paramètres physiques de l'environnement atomique en un seul
cycle expérimental, ainsi que deux observables non commutantes d'un qubit en une seule
réalisation, en transférant son état sur des états de spin auxiliaires. Cette dernière
démonstration est décrite comme une POVM (\textit{positive operator-valued measure}).

Nous étudions ensuite la photoassociation près de la raie $^1S_0\rightarrow{}^3P_1$. La
spectroscopie que nous avons réalisée sur l'isotope bosonique $^{88}$Sr a servi à valider
les méthodes extrayant la longueur optique à partir des pertes atomiques. Pour le $^{87}$Sr,
la théorie prévoit jusqu'à quelques milliers de raies étroites dans la plage explorée de
$400$~MHz. Nous n'observons que trois raies, élargies à plusieurs MHz, soit deux à trois
ordres de grandeur au-dessus de leur largeur naturelle, avec des longueurs optiques
inférieures d'un à deux ordres de grandeur à celles du $^{88}$Sr. Nous discutons enfin
plusieurs interprétations possibles de ces résultats, les plus probables étant le mécanisme
de prédissociation, la statistique quantique et les effets de sélection de spin.

\vspace{0.15cm}
\noindent\textbf{Mots clefs :} strontium 87, spin nucléaire, qudits, symétrie SU(N), interférométrie de Ramsey, POVM, photoassociation, gaz de Fermi dégénéré.
\par}
